\section{Introduction}
\IEEEPARstart{L}{ow} Earth orbit (LEO) constellations such as Starlink and OneWeb now give remote and airborne terminals broadband service where no terrestrial network reaches, and 3GPP integrates such non-terrestrial links into 5G~\cite{lin20215g}. A terminal that has both a LEO and a terrestrial path can send each packet on the path that will deliver it in time. Multipath transport protocols such as Multipath QUIC (MPQUIC) provide the mechanism~\cite{deconinck2017multipath}. Their schedulers react to smoothed RTT estimates of the recent past~\cite{paasch2014experimental,lim2017ecf}. On LEO paths this is too slow. Starlink latency jumps at globally synchronized reconfiguration instants every 15\,s~\cite{mohan2024multifaceted,tanveer2023making}. OneWeb latency swings at satellite handovers~\cite{zhao2025measuring}. Measured MPTCP and MPQUIC sessions over Starlink and cellular paths show scheduler failures caused by these dynamics~\cite{baltaci2023multipath}. Schedulers that use forecasts of path state instead of measurements of it have been proposed for multi-access and mobile settings~\cite{yang2024mams,maglione2025intelligent}. They need forecasts that are accurate a few hundred milliseconds to a few seconds ahead and, because a wrong forecast costs a missed deadline, they need to know how wrong each forecast may be.

Prior work on LEO performance prediction targets throughput or latency at coarse time scales or over large areas~\cite{liu2025vivisecting,benghe2026horizon}. Attention models have also been used to predict LEO RTT~\cite{tian2025attention}. Work on terrestrial links shows that fine-grained throughput is hard to predict and that learned models should be compared with simple baselines~\cite{basit2024predictability,raca2020leveraging}. Three questions remain open for the scheduling use case. Are the measurement pipelines that produce the training data causal? Which point forecast should be reported for heavy-tailed link metrics, where a few spikes dominate the squared error? And can the uncertainty of a sub-second forecast be calibrated on sessions the model has not seen?

This paper answers these questions with measurements of a Starlink, a OneWeb, and a 5G link taken in parallel at 5\,Hz on the ARA wireless living lab~\cite{islam2025design}. Its contributions are the following.
\begin{enumerate}
  \item \textbf{A causality audit of link-state data.} We show that aligning ping streams by time interpolation, a common step when several collectors are merged on one time base, places part of the next sample in the model input. We prove that the leak reduces the one-step persistence error of every interpolated row by the interpolation weight (one half for evenly spaced rows) and measure that it hides \AuditPersHonePct\% of the first-step RTT error on Starlink downlink, where \InterpFracSLDL\% of the samples are interpolated. All results in this paper use causal inputs.
  \item \textbf{A distributional encoder-decoder.} We extend a GRU sequence-to-sequence forecaster with additive attention~\cite{cho2014learning,bahdanau2015neural} by anchoring its decoder on the last observation and replacing each scalar head with a head that outputs a mean and five quantiles whose order is guaranteed by construction. A single loss trains the mean and the quantiles, and each output is the optimal forecast under its own scoring rule~\cite{gneiting2011making}.
  \item \textbf{Calibration on unseen sessions.} We calibrate the quantile intervals with conformalized quantile regression (CQR)~\cite{romano2019conformalized} on held-out measurement sessions and report the coverage measured on further held-out sessions, instead of coverage computed from the test residuals.
  \item \textbf{A controlled evaluation.} On six links, with three seeds, an ablation of each design step, and four non-neural references, the proposed model improves the 95\% interval score of the RTT forecast by \ChgMISWorst\% to \ChgMISBest\% and its mean absolute error by \ChgMAEWorst\% to \ChgMAEBest\% against the same network with a scalar output, while the root-mean-square error stays within a few percent. We report where it does not help, including a spike-warning task on OneWeb downlink.
\end{enumerate}

Section~\ref{sec:related} reviews related work. Section~\ref{sec:system} describes the measurement system and formulates the forecasting problem, including the causality analysis. Section~\ref{sec:method} presents the forecaster and its calibration. Sections~\ref{sec:setup} and~\ref{sec:results} give the evaluation protocol and the results, and Section~\ref{sec:discussion} discusses their use for multipath scheduling and concludes the paper.
